\cleardoublepage

\chapter{Pamięć podręczna i~serwer pośredniczący na~ścieżce żądania}

\label{cha:WzorceProjektowe}

Na~ścieżce żądania w~badanym systemie występują dwa elementy wpływające na~zmierzony czas operacji niezależnie od~wybranego protokołu: pamięć podręczna po~stronie usługi oraz serwer pośredniczący w~jednym z~wariantów gRPC-Web. Oba wymagają omówienia przed przejściem do~części badawczej.

% -----------------------------------------------------------------------
\section{Wzorzec Cache-Aside}
\label{sec:CacheAside}
% -----------------------------------------------------------------------

Wzorzec \textit{Cache-Aside} polega na~bezpośrednim zarządzaniu pamięcią podręczną przez warstwę aplikacji, która odpowiada za~synchronizację pomiędzy magazynem podręcznym a~głównym źródłem danych~\cite{bib:CacheAside}. W~trakcie operacji odczytu aplikacja weryfikuje dostępność żądanego zasobu w~magazynie podręcznym na~podstawie deterministycznie wygenerowanego klucza. Przy trafieniu (ang.\ \textit{cache hit}) dane zwracane są~z~magazynu podręcznego, bez odpytywania relacyjnej bazy danych. Przy chybieniu (ang.\ \textit{cache miss}) aplikacja kieruje zapytanie do~głównej bazy danych, zapisuje otrzymany wynik w~pamięci podręcznej z~określonym czasem ważności (TTL) i~dopiero wtedy odsyła go~do~klienta, przez co~bufor gromadzi wyłącznie dane rzeczywiście odpytywane przez klientów. Nieaktualne wpisy unieważnia się (ang.\ \textit{cache invalidation}) pasywnym wygasaniem czasowym TTL albo jawnym usunięciem klucza w~momencie zapisu danych w~głównym magazynie.

W~badanym systemie rolę magazynu podręcznego pełni Redis, uruchomiony jako odrębna usługa. Trafienie eliminuje ze~ścieżki żądania zapytanie do~relacyjnej bazy danych, wymaga jednak wymiany komunikatów z~pamięcią podręczną.

% -----------------------------------------------------------------------
\section{Rola serwera pośredniczącego Envoy}
\label{sec:Envoy}
% -----------------------------------------------------------------------

W~architekturach zorientowanych na~usługi serwery pośredniczące (ang.\ \textit{reverse proxy} lub \textit{API Gateway}) obsługują ruch przychodzący, agregują żądania oraz tłumaczą pomiędzy protokołami komunikacyjnymi. Envoy Proxy, napisany w~języku C++, udostępnia w~tym celu filtr \texttt{grpc\_web}, który tłumaczy ruch przeglądarkowy gRPC-Web na natywne gRPC oparte na HTTP/2~\cite{bib:EnvoyProxy}. Usługi biznesowe nie muszą wówczas implementować translacji protokołów ani reguł kierowania ruchem. Na~ścieżkę żądania wchodzi dodatkowy przeskok sieciowy (ang.\ \textit{network hop}) oraz koszt przetwarzania i~buforowania ramek danych.

W~badanym systemie Envoy realizuje jeden z~dwóch porównywanych wariantów gRPC-Web, drugi opiera się na~obsłudze bezpośrednio w~usłudze ASP.NET Core. Obecność proxy należy uwzględniać przy interpretacji wyników, ponieważ różnica między wariantem REST a~ścieżką z~serwerem pośredniczącym obejmuje jednocześnie sposób kodowania komunikatów, translację gRPC-Web oraz dodatkowy element pośredniczący. Dlatego w~analizie traktowano Envoy jako część całej konfiguracji komunikacyjnej, a~porównania wykonywano dla pełnych ścieżek żądania: od~przeglądarki do~usługi i~z~powrotem.
